% Ported verbatim from paper/top3_joint_moment_bound.tex lines 718-816.
% Only cross-reference wording was adapted; no mathematics was altered.
We now turn to the numerical estimate. Throughout this section the law
is arbitrary and isotropic. Bounded positive radii will be imposed only
when a radial comparison requires them. For later bookkeeping define
\begin{equation}\label{constants}
 d_0=-\frac{25624284539519}{53655078400000}
                  +\frac{\sqrt5}{2}-\frac{41\sqrt3}{120},\qquad
 d_1=d_0+\frac{103}{9000000},\qquad
 d=d_1+\frac1{3000000}.
\end{equation}
Then $d=(3+\sqrt5)/2-2c$, with $c$ as in~\eqref{main-constant}.
The intermediate constants separate the joint moment estimate from
the final atomic term.
Fix $\E XX^T=I_3$, let $R=|X|$, and put $U=X/R$ when $R>0$.
At $R=0$ choose any fixed unit vector for $U$. Define
\[
 \begin{gathered}
 C=\E R,\quad V=3-C^2,\quad \mathbf m=\E X,\quad B=|\mathbf m|^2,\quad e=1-B,\\
 A=\E|X\cdot X'|,\qquad D=\frac{3+\sqrt5}{2}-A-B.
 \end{gathered}
\]
Here $V,e\ge0$: $\E R^2=3$ and Cauchy--Schwarz gives
$(v\cdot \mathbf m)^2\le\E(v\cdot X)^2=1$ for every unit $v$.
All pair expectations below are over independent copies of this same law.
Set
\[
 M=\E[R UU^T]=\frac C3I_3+Z,\qquad
 J=1-\tr(M^2)=\frac V3-\|Z\|_F^2.
\]
The regression vector $\xi=U-MX$ satisfies
\begin{equation}\label{regression}
 \E[X\xi^T]=0,\quad \E\|\xi\|^2=J\ge0,\quad
 \E UU^T=M^2+\E[\xi\xi^T].
\end{equation}
In particular $J$ is an actual regression error, not a freely chosen moment.

Here are the harmonic normalizations used throughout. All tensor inner
products sum over all ordered indices. Put
\[
 \Omega_{ijkl}=\frac{\delta_{ij}\delta_{kl}+\delta_{ik}\delta_{jl}
                           +\delta_{il}\delta_{jk}}{15},
 \qquad
 ({\cal L}_3(y))_{ijk}
       =\frac{y_i\delta_{jk}+y_j\delta_{ik}+y_k\delta_{ij}}5.
\]
For symmetric traceless $Z$, put
\[
 ({\cal L}_4(Z))_{ijkl}=\frac17(\delta_{ij}Z_{kl}+\delta_{ik}Z_{jl}
 +\delta_{il}Z_{jk}+\delta_{jk}Z_{il}+\delta_{jl}Z_{ik}+\delta_{kl}Z_{ij}).
\]
For unit $u$, set
\[
 \begin{gathered}
 \Psi_3(u)=u^{\otimes3}-{\cal L}_3(u),\qquad
 \Psi_4(u)=u^{\otimes4}-\Omega-{\cal L}_4(uu^T-I_3/3),\\
 T_i=\E[R\Psi_i(U)],\qquad H=\|T_3\|_F,\qquad K=\|T_4\|_F.
 \end{gathered}
\]
Contraction in two indices gives the orthogonal decompositions
\[
 \E[R U^{\otimes3}]=T_3+{\cal L}_3(\mathbf m),\qquad
 \E[R U^{\otimes4}]=T_4+C\Omega+{\cal L}_4(Z).
\]
Indeed $\operatorname{ct}\Omega=I_3/3$, $\|\Omega\|^2=1/5$,
$\operatorname{ct}{\cal L}_4(Z)=Z$, and
$\langle T,{\cal L}_4(Z)\rangle=(6/7)\langle\operatorname{ct}T,Z\rangle$.
Also $\operatorname{ct}{\cal L}_3(y)=y$ and
$\langle N,{\cal L}_3(y)\rangle=(3/5)\operatorname{ct}N\cdot y$.
It follows directly that
\begin{align}\label{harmonics}
 \langle\Psi_3(u),\Psi_3(v)\rangle&=g_3(u\cdot v),&
 g_3(t)&=t^3-3t/5,&\|\Psi_3(v)\|^2&=2/5,\notag\\
 \langle\Psi_4(u),\Psi_4(v)\rangle&=g_4(u\cdot v),&
 g_4(t)&=t^4-6t^2/7+3/35,&\|\Psi_4(v)\|^2&=8/35.
\end{align}
The tensor moments involve $R$, not $R^3$ or $R^4$, and are integrable
under $\E R^2=3$.

Write $s=\sqrt5$ and define the scalar absolute-value majorant and its slack by
\[
 a=\frac{1+s}{6},\quad b=\frac{7s-5}{8},\quad\rho=\frac{15s-25}{16},
 \qquad h_0(z)=\frac{(1-z)(sz-1)^2(sz+s+2)}{2(s+1)^2}.
\]
Expansion gives, for $|t|\le1$,
\[
 a+b(t^2-1/3)-\rho(t^4-1/5)-|t|=h_0(|t|)\ge0.
\]
This majorant is from~\cite[Lemma 4.4]{STW}.
Let $S=\E[RR'h_0(|U\cdot U'|)]$. The decompositions above give
\[
 A=aC^2+b\|Z\|^2-\rho\left(K^2+\frac67\|Z\|^2\right)-S.
\]
Consequently
\begin{equation}\label{budget4}
 D=e+\alpha V+\beta J+\rho K^2+S,
 \qquad \alpha=\frac{2s-1}{14},\quad\beta=\frac{s+10}{14}.
\end{equation}
Useful rational bounds are $\alpha\ge26/105$, $\beta\ge6/7$,
and $\rho\ge1/2$, using $s>67/30$, $s\ge2$, and $s\ge11/5$.
